\documentclass[12pt,a4paper]{article}

% ============================================================
% ЯЗЫКИ И ШРИФТЫ (XeLaTeX)
% ============================================================
\usepackage{fontspec}
\usepackage{polyglossia}
\setdefaultlanguage{russian}
\setotherlanguage{english}
\setmainfont{CMU Serif}
\setsansfont{CMU Sans Serif}
\setmonofont{CMU Typewriter Text}

\usepackage{geometry}
\geometry{top=2cm, bottom=2cm, left=2.5cm, right=2.5cm}

\usepackage{amsmath,amssymb,amsthm}
\usepackage{booktabs}
\usepackage{longtable}
\usepackage{float}
\usepackage{xcolor}
\usepackage{hyperref}
\hypersetup{colorlinks, urlcolor=blue, citecolor=blue}

\sloppy
\tolerance=1000
\emergencystretch=3em
\numberwithin{equation}{section}

\begin{document}

\begin{center}
    {\LARGE\bf Модель MEML 75}\\[0.3cm]
    {\large Вывод калибровочной константы, фактора подавления\\
    и константы распада аксиона из геометрии орбифолда $T^2/\mathbb{Z}_3$}\\[0.5cm]
    {\bf Юров Д. В.}\\[0.2cm]
    {\it 2026}
\end{center}

\vspace{0.5cm}

\begin{abstract}
Версия MEML~75 развивает аксионную парадигму на орбифолде 
$T^2/\mathbb{Z}_3$ в шестимерном пространстве-времени. Bulk-константа 
связи $\alpha_G = 2/(3\pi) = 0.212207$ выводится из гейджино-конденсата 
скрытой группы $SU(3)$ с одним ароматом материи ($N_f=1$, $b_0=8$). 
Фактор подавления $\varepsilon = 2.860924\times 10^{-4}$ выводится 
из суммирования по спектру Калуцы--Клейна. Используется стандартная 
конвенция приведённой массы Планка $M_{Pl}^{\text{red}} = 
2.435\times 10^{18}$~ГэВ, что даёт масштаб компактификации 
$m_{KK} = 2.084\times 10^{14}$~ГэВ и 6D фундаментальный масштаб 
$M_6 = 2.965\times 10^{16}$~ГэВ, естественным образом совпадающий 
с масштабом Великого Объединения. \textbf{Ключевой результат:} 
константа распада аксиона выводится из первых принципов как 
$f_a = M_{Pl}^{\text{red}}\varepsilon^2 = 1.993\times 10^{11}$~ГэВ, 
без каких-либо феноменологических нормировок. Это даёт массу аксиона 
$m_a = 28.22$~мкэВ и резонансную частоту $\nu_a = 6.825$~ГГц, 
попадающую в диапазон высокочастотных галоскопов ADMX-HF / HAYSTAC / 
CAPP. Связь аксиона с фотонами составляет 
$|g_{a\gamma\gamma}| = 1.119\times 10^{-14}$~ГэВ$^{-1}$. 
Плотность тёмной материи $\Omega_a h^2 = 0.119832$ (Planck 2018: 
$0.120 \pm 0.001$) воспроизводится при начальном угле смещения 
$\theta_i = 2.0043$. Начальный угол $\theta_i$ является 
\textbf{начальным условием} космологической эволюции аксиона и не 
является параметром теории. Единственный калибровочный параметр 
модели~---~глубина варп-ямы $k = 0.1068$.
\end{abstract}

\section{Конвенции}
\label{sec:conventions}

\subsection{Масса Планка}
\begin{equation}
M_{Pl}^{\text{red}} = \frac{1}{\sqrt{8\pi G}} 
= 2.435\times 10^{18}~\text{ГэВ}.
\end{equation}

\subsection{6D действие}
\begin{equation}
S_6 = \frac{M_6^4}{2}\int d^6x\sqrt{-G}\,R_6.
\end{equation}
После редукции на $M_4\times T^2/\mathbb{Z}_3$:
\begin{equation}
(M_{Pl}^{\text{red}})^2 = \frac{M_6^4}{3m_{KK}^2}.
\end{equation}

\subsection{Аксионный сектор}
\begin{equation}
\boxed{\;f_a = M_{Pl}^{\text{red}}\,\varepsilon^2 = 
1.993\times 10^{11}~\text{ГэВ}.\;}
\end{equation}

\section{Введение}
Версия MEML~10 содержала три свободных параметра ($\alpha_G$, 
$\varepsilon$, $\theta_i$). Ревизия~75 устраняет первые два, выводя 
их из геометрии компактифицированного пространства.

\section{Геометрия орбифолда $T^2/\mathbb{Z}_3$}

\subsection{Комплексная структура}
\begin{equation}
\tau = e^{i\pi/3} = \frac{1}{2} + i\frac{\sqrt3}{2}, \qquad 
z = x + \tau y, \quad x,y\in[0,1).
\end{equation}

\subsection{Метрика тора}
\begin{equation}
ds^2 = dx^2 + dx\,dy + dy^2.
\end{equation}

\subsection{Действие $\mathbb{Z}_3$}
\begin{equation}
(x,y) \to (-x-y,\, x).
\end{equation}

\subsection{Неподвижные точки}
\begin{equation}
z_1 = (0,0), \quad z_2 = \left(\tfrac13,\tfrac13\right), 
\quad z_3 = \left(\tfrac23,\tfrac23\right), \qquad R^2 = 1/3.
\end{equation}

\subsection{Площадь фундаментальной области}
\begin{equation}
V_{\text{orb}} = \frac{V_2}{3} = \frac{1}{3m_{KK}^2}.
\end{equation}

\section{Связь масштабов}
\begin{equation}
M_6 = \left(3(M_{Pl}^{\text{red}})^2 m_{KK}^2\right)^{1/4} 
= 2.965\times 10^{16}~\text{ГэВ}.
\end{equation}
Совпадает с $M_{GUT} \sim 10^{16}$~ГэВ.

\section{Вывод $\alpha_G$}
\begin{equation}
\alpha_G = \frac{b_0}{12\pi}, \qquad b_0 = 3N - N_f.
\end{equation}

\begin{table}[h!]
\centering
\caption{Скан скрытых групп.}
\begin{tabular}{|c|c|c|c|c|}
\hline
Группа & $N_f$ & $b_0$ & $\alpha_G$ & Отклонение \\
\hline
$SU(2)$ & 0 & 6 & 0.159155 & $-24.67\%$ \\
$SU(3)$ & 0 & 9 & 0.238732 & $+12.99\%$ \\
$\mathbf{SU(3)}$ & $\mathbf{1}$ & $\mathbf{8}$ & $\mathbf{0.212207}$ 
& $\mathbf{+0.44\%}$ \\
$SU(3)$ & 2 & 7 & 0.185681 & $-12.12\%$ \\
$SU(4)$ & 4 & 8 & 0.212207 & $+0.44\%$ \\
\hline
\end{tabular}
\end{table}

\begin{equation}
\alpha_G = \frac{2}{3\pi} = 0.21220659, \qquad 
m_{KK} = \frac{M_{Pl}^{\text{red}}}{\sqrt{\mathcal{I}}\,e^{\pi^2}} 
= 2.084\times 10^{14}~\text{ГэВ}.
\end{equation}

\section{Спектр Калуцы--Клейна}
\begin{equation}
m^2(m,n) = \frac{2}{\sqrt3}\,m_{KK}^2\,(m^2 - mn + n^2).
\end{equation}

\section{Вывод $\varepsilon$}
\begin{equation}
x = \frac{\alpha_G}{4\pi} = \frac{1}{6\pi^2} = 0.0168868639.
\end{equation}
\begin{equation}
\varepsilon = x^2\left[1 + c_3'\,x + c_4'\,x^2 + c_5'\,x^3\right], 
\qquad c_3' = 0.192, \; c_4' = 0.0203, \; c_5' = 0.00271.
\end{equation}
\begin{equation}
\varepsilon = 2.860924\times 10^{-4}.
\end{equation}

\section{Модулярный интеграл $\mathcal{I}$}
\begin{equation}
\mathcal{I} = \frac{V_{\text{orb}}^2}
{\int_{T^2/\mathbb{Z}_3} e^{2A}\,d^2z} = \frac{V_{\text{orb}}}{J}.
\end{equation}
Монте-Карло ($N = 10^8$):
\begin{equation}
J = 0.912553 \pm 0.000004, \qquad \mathcal{I} = 0.365275 \pm 0.000002.
\end{equation}

\section{Вывод $f_a$ из первых принципов}
\begin{equation}
\boxed{\;f_a = M_{Pl}^{\text{red}}\,\varepsilon^2 
= 1.993\times 10^{11}~\text{ГэВ}.\;}
\end{equation}
\textbf{Ключевое отличие от MEML~10:} $f_a$ не подгоняется, 
а выводится из первых принципов.

\section{Космологические следствия}
\begin{align}
m_a &= \frac{\Lambda_{\text{QCD}}^2}{f_a} = 28.22~\text{мкэВ},\\
\nu_a &= \frac{m_a c^2}{h} = 6.825~\text{ГГц},\\
\Omega_a h^2 &= 0.119832 \quad (\text{Planck 2018: } 0.120 \pm 0.001),\\
|g_{a\gamma\gamma}| &= 1.119\times 10^{-14}~\text{ГэВ}^{-1}.
\end{align}

\section{Астрофизический модуль: ПЧД}
Для суперрадиационной неустойчивости моды $l=m=1$:
\begin{equation}
M_{BH} = \frac{\alpha_{\max}\,\hbar c}{G\,m_a} 
= 1.894\times 10^{-6}\,M_\odot \approx 0.63\,M_\oplus.
\end{equation}
\begin{equation}
T_H = \frac{\hbar c^3}{8\pi G M_{BH} k_B} = 32.0~\text{мкэВ} 
> m_a = 28.2~\text{мкэВ}.
\end{equation}
\begin{equation}
\tau_{SR} = 0.85~\mu\text{с}.
\end{equation}

\section{Статус параметров}
\begin{table}[h!]
\centering
\caption{Статус независимых параметров MEML~75.}
\begin{tabular}{|l|c|l|l|}
\hline
Величина & Значение & Статус & Источник \\
\hline
$\alpha_G$ & $0.212207$ & Выведен & $SU(3)$, $N_f=1$ \\
$m_{KK}$ & $2.084\times 10^{14}$~ГэВ & Выведен & Радион \\
$\varepsilon$ & $2.861\times 10^{-4}$ & Выведен & KK-сумма \\
$f_a$ & $1.993\times 10^{11}$~ГэВ & \textbf{Выведен} & 
$M_{Pl}^{\text{red}}\varepsilon^2$ \\
$\mathcal{I}$ & $0.365275$ & Условно выведен & Гравитон + $k$ \\
$k$ & $0.1068$ & Калиброван & Под $\mathcal{I}$ \\
$\sigma_0^2/V_{\text{orb}}$ & $0.06$ & Модельный & Оценка \\
$\theta_i$ & $2.0043$ & Начальное условие & Космология \\
\hline
$M_6$ & $2.965\times 10^{16}$~ГэВ & Производная & 
$(3(M_{Pl}^{\text{red}})^2 m_{KK}^2)^{1/4}$ \\
\hline
\end{tabular}
\end{table}

\section{Заключение}
MEML~75 — частично предсказательная модель с $f_a$, впервые выведенной 
из первых принципов. Единственный калибровочный параметр~---~$k = 0.1068$. 
Ключевые предсказания: $\nu_a = 6.825$~ГГц (ADMX-HF), 
$|g_{a\gamma\gamma}| = 1.119\times 10^{-14}$~ГэВ$^{-1}$, 
$M_{BH} \approx 0.63\,M_\oplus$ (ПЧД с открытым аксионным каналом Хокинга).

\begin{thebibliography}{9}
\bibitem{Planck2018} Planck Collaboration, Astron. Astrophys. \textbf{641}, A6 (2020).
\bibitem{PecceiQuinn1977} R.~D. Peccei, H.~R. Quinn, Phys. Rev. Lett. \textbf{38}, 1440 (1977).
\bibitem{LythTurner} D.~H. Lyth, M.~S. Turner, Phys. Rev. D \textbf{48}, 1502 (1993).
\bibitem{ADMX2025} ADMX Collaboration, Phys. Rev. Lett. \textbf{134}, 021002 (2025).
\bibitem{KKLT2003} S. Kachru, R. Kallosh, A. Linde, S. Trivedi, Phys. Rev. D \textbf{68}, 046005 (2003).
\bibitem{Sikivie1982} P.~Sikivie, Phys. Rev. Lett. \textbf{48}, 1156 (1982).
\bibitem{HAYSTAC2023} HAYSTAC Collaboration, Phys. Rev. Lett. \textbf{130}, 201001 (2023).
\end{thebibliography}

\end{document}